% v2.4.0 P2 figure batch 02: V5-C01 plus the first V5-C02 P2 UID.
\documentclass[UTF8,a4paper,11pt,openany]{ctexbook}
\usepackage{../common/statlearnbook}
\SLSetBookMetadata{P2 绘图批次 02}{统计学习方法讲义 v2.4.0 绘图集成验证}{第五册第 30 章与第 31 章首图局部编译}
\begin{document}
\pagestyle{empty}
\raggedbottom

\typeout{SLFIG-HARNESS-BEGIN FIG-P544-01}
% v2.6.0 figure UID: FIG-P544-01
\tikzset{
  slfig-FIG-P544-01/.style={font=\fontsize{9.4pt}{11.2pt}\selectfont,
    every path/.append style={line cap=round,line join=round}},
  slfig-FIG-P544-01-base/.style={draw=SLBlue,rounded corners=2pt,fill=SLBlue!7,
    minimum width=39mm,minimum height=9mm,line width=.84pt,align=center},
  slfig-FIG-P544-01-mid/.style={draw=SLTeal,rounded corners=2pt,fill=SLTeal!7,
    minimum width=39mm,minimum height=10mm,line width=.84pt,align=center},
  slfig-FIG-P544-01-top/.style={draw=SLGold,rounded corners=2pt,fill=SLGold!10,
    minimum width=54mm,minimum height=10mm,line width=.92pt,align=center},
  slfig-FIG-P544-01-reversible/.style={draw=SLGold,rounded corners=2pt,fill=white,
    minimum width=31mm,minimum height=9mm,line width=.78pt,align=center},
  slfig-FIG-P544-01-dep/.style={-{Stealth[length=2.2mm]},draw=SLBlue,line width=1.0pt},
  slfig-FIG-P544-01-sufficient/.style={-{Stealth[length=2mm]},draw=SLGold,
    line width=.72pt,densely dashed},
  slfig-FIG-P544-01-edgelabel/.style={font=\fontsize{8.8pt}{10.4pt}\selectfont,
    text=SLGold,fill=white,inner sep=.8pt},
  slfig-FIG-P544-01-legend/.style={anchor=west,text=SLTextGray,
    font=\fontsize{8.8pt}{10.4pt}\selectfont,inner sep=0pt}
}
\begin{figure}[htbp]
\centering
\begin{tikzpicture}[slfig-FIG-P544-01,>=Stealth]
  \node[slfig-FIG-P544-01-base] (foundation) at (0,-2.05) {马尔可夫性质 / 转移规则};
  \node[slfig-FIG-P544-01-mid] (fixed) at (-2.55,-.15) {平稳固定点\\$\pi=\pi P$};
  \node[slfig-FIG-P544-01-mid] (structure) at (2.55,-.15) {不可约 / 非周期 / 返性};
  \node[slfig-FIG-P544-01-top] (longrun) at (0,1.85) {长期结论\\时间平均与逐步收敛};
  \node[slfig-FIG-P544-01-reversible] (reversible) at (-5.25,-2.05)
    {可逆性 / 细致平衡};

  \draw[slfig-FIG-P544-01-dep] (foundation)--(fixed);
  \draw[slfig-FIG-P544-01-dep] (foundation)--(structure);
  \draw[slfig-FIG-P544-01-dep] (fixed)--(longrun);
  \draw[slfig-FIG-P544-01-dep] (structure)--(longrun);
  \draw[slfig-FIG-P544-01-sufficient] (reversible)--(fixed)
    node[slfig-FIG-P544-01-edgelabel,midway,above left] {充分条件};

  \node[slfig-FIG-P544-01-legend] at (-5.75,2.65)
    {实线：必要依赖链};
  \draw[slfig-FIG-P544-01-dep,yshift=1.8mm] (-3.65,2.66)--(-2.75,2.66);
  \node[slfig-FIG-P544-01-legend] at (-1.95,2.65)
    {虚线：仅充分，不是必要前置};
  \draw[slfig-FIG-P544-01-sufficient,yshift=1.8mm] (1.20,2.66)--(2.10,2.66);
\end{tikzpicture}
\caption{本章概念依赖：从马尔可夫性质和转移规则出发，分别经平稳固定点与结构条件到达长期结论；可逆性只提供通向平稳性的充分条件}\label{fig:V5-C01-dependency}
\end{figure}

\clearpage
\typeout{SLFIG-HARNESS-END FIG-P544-01}

\typeout{SLFIG-HARNESS-BEGIN FIG-P547-01}
\input{../../绘图源码/第05册_采样方法主题模型与图排序/V5-C01/fig_v5_c01_transition_graph.tex}
\clearpage
\typeout{SLFIG-HARNESS-END FIG-P547-01}

\typeout{SLFIG-HARNESS-BEGIN FIG-P552-01}
% v2.3.1 figure UID: FIG-P552-01
% v2.6.0 reverse redraw for FIG-P552-01: shared timeline grid and explicit readable hierarchy.
\tikzset{slfig-FIG-P552-01/.style={font=\fontsize{9.4pt}{11.2pt}\selectfont,every path/.append style={line cap=round,line join=round}}}
\begin{figure}[htbp]
\centering
\begin{tikzpicture}[slfig-FIG-P552-01,x=1.18cm,y=1cm,>=Stealth,
  event/.style={circle,draw=SLGray,fill=SLBlue!5,minimum size=6mm,line width=.68pt,
    font=\fontsize{9.2pt}{11.0pt}\selectfont},
  hit/.style={event,draw=SLGold,fill=SLGold!14,line width=.92pt},
  timeline/.style={-{Stealth[length=1.9mm]},draw=SLTextGray,line width=.78pt},
  interval/.style={-{Stealth[length=2.1mm]},draw=SLGold,line width=1pt},
  interval label/.style={font=\fontsize{9.2pt}{11.0pt}\selectfont,fill=white,inner sep=1.1pt},
  tick label/.style={font=\fontsize{8.6pt}{10.1pt}\selectfont},
  tick mark/.style={draw=SLTextGray,line width=.5pt},
  title/.style={font=\fontsize{9.4pt}{11.2pt}\selectfont\bfseries},
  note/.style={font=\fontsize{8.8pt}{10.4pt}\selectfont,fill=white,inner sep=1.1pt},
  conclusion/.style={draw=SLRuleGray,rounded corners=2pt,fill=white,align=center,
    minimum width=30mm,minimum height=12mm,line width=.68pt,inner sep=2.2pt,
    font=\fontsize{9.2pt}{11.0pt}\selectfont},
]
  \node[anchor=east,title,text=SLBlue] at (-.35,1.25) {首次到达 $\tau_i$};
  \draw[timeline] (0,1.25)--(7.45,1.25) node[right] {$t$};
  \foreach \t/\s in {0/j,1/k,2/k,3/i,4/k,5/i,6/k,7/k}{
    \draw[tick mark] (\t,1.33)--(\t,1.17);
    \node[below=2mm,tick label] at (\t,1.17) {$\t$};
    \node[event] at (\t,1.78) {$\s$};
  }
  \node[hit] at (3,1.78) {$i$};
  \draw[interval] (0,2.30)--(3,2.30)
    node[midway,above,interval label] {$j\ne i$ 出发，首次到达：$\tau_i=3$};

  \node[anchor=east,title,text=SLTeal] at (-.35,-1.10) {首次正回返 $\tau_i^+$};
  \draw[timeline] (0,-1.10)--(7.45,-1.10) node[right] {$t$};
  \foreach \t/\s in {0/i,1/k,2/k,3/i,4/k,5/i,6/k,7/k}{
    \draw[tick mark] (\t,-1.02)--(\t,-1.18);
    \node[below=2mm,tick label] at (\t,-1.18) {$\t$};
    \node[event] at (\t,-.57) {$\s$};
  }
  \node[hit] at (3,-.57) {$i$};
  \node[anchor=west,note,text=SLGold] at (.15,-.05) {$t>0$ 才计入};
  \draw[interval] (0,.10)--(3,.10)
    node[midway,above,interval label] {从 $i$ 出发，首次正回返：$\tau_i^+=3$};

  \node[conclusion] at (8.85,.10)
    {正常返考察\\$\mathbb E_i[\tau_i^+]<\infty$};
\end{tikzpicture}
\caption{首次到达与首次正回返的区分：正常返考察从$i$出发的$\mathbb E_i[\tau_i^+]$，而不是某个固定时刻首次到达的概率}\label{fig:V5-C01-return-time}
\end{figure}

\clearpage
\typeout{SLFIG-HARNESS-END FIG-P552-01}

\typeout{SLFIG-HARNESS-BEGIN FIG-P556-01}
% v2.3.1 figure UID: FIG-P556-01
% Proposed post-v2.3.1 polish for FIG-P556-01: protect key-label size and stroke hierarchy.
\tikzset{slfig-FIG-P556-01/.style={every path/.append style={line cap=round,line join=round}}}
\pgfplotsset{
  slfig-FIG-P556-01-axis/.style={
    tick label style={font=\fontsize{8.7pt}{10.2pt}\selectfont},
    label style={font=\fontsize{9.4pt}{11.2pt}\selectfont},
    clip=false,
  },
  slfig-FIG-P556-01-reference/.style={draw=SLGray,line width=.58pt,densely dashed},
  slfig-FIG-P556-01-map/.style={draw=SLTeal,line width=1.02pt},
  slfig-FIG-P556-01-cobweb-low/.style={draw=SLBlue,line width=.72pt},
  slfig-FIG-P556-01-cobweb-high/.style={draw=SLGold,line width=.72pt,densely dashed},
}
\begin{figure}[htbp]
\centering
\begin{tikzpicture}[slfig-FIG-P556-01]
\begin{axis}[slfig-FIG-P556-01-axis,
  width=9.2cm,height=7.3cm,
  xmin=0,xmax=1,ymin=0,ymax=1,
  axis equal image,axis lines=left,
  xlabel={$r_t$},ylabel={$r_{t+1}$},
  xtick={0,.2,.4,.6,.8,1},ytick={0,.2,.4,.6,.8,1},
  clip=false,
]
  \addplot[slfig-FIG-P556-01-reference,domain=0:1,samples=201] {x}
    node[pos=.88,above left,font=\fontsize{9.3pt}{11.0pt}\selectfont] {$y=r$};
  \addplot[slfig-FIG-P556-01-map,domain=0:1,samples=201] {.2+.5*x}
    node[pos=.94,above,yshift=1.6mm,font=\fontsize{9.3pt}{11.0pt}\selectfont,text=SLTeal] {$y=0.2+0.5r$};
  \addplot[only marks,mark=*,mark size=2.7pt,draw=SLGold,fill=SLGold]
    coordinates {(.4,.4)} node[above left,xshift=-.8mm,yshift=.8mm,font=\fontsize{9.4pt}{11.2pt}\selectfont,text=SLGold] {$r^*=0.4$};

  % r_0=0.05：实线 cobweb。
  \addplot[slfig-FIG-P556-01-cobweb-low] coordinates {
    (.05,0)(.05,.225)(.225,.225)(.225,.3125)(.3125,.3125)
    (.3125,.35625)(.35625,.35625)(.35625,.378125)(.378125,.378125)};
  \addplot[only marks,mark=square*,mark size=2.1pt,draw=SLBlue,fill=SLBlue]
    coordinates {(.05,0)};
  \node[anchor=west,yshift=-.6mm,font=\fontsize{9.2pt}{11.0pt}\selectfont,text=SLBlue] at (axis cs:.08,.075) {$r_0=.05$};

  % r_0=0.90：虚线 cobweb。
  \addplot[slfig-FIG-P556-01-cobweb-high] coordinates {
    (.90,0)(.90,.65)(.65,.65)(.65,.525)(.525,.525)
    (.525,.4625)(.4625,.4625)(.4625,.43125)(.43125,.43125)};
  \addplot[only marks,mark=triangle*,mark size=2.5pt,draw=SLGold,fill=SLGold]
    coordinates {(.90,0)};
  \node[anchor=east,yshift=-1.2mm,font=\fontsize{9.2pt}{11.0pt}\selectfont,text=SLGold] at (axis cs:.87,.075) {$r_0=.90$};
  \node[anchor=east,draw=SLRuleGray,line width=.65pt,rounded corners=2pt,fill=white,inner sep=2.2pt,align=center,font=\fontsize{9.2pt}{11.0pt}\selectfont]
    at (axis cs:.74,.18) {斜率 $0.5<1$\\压缩映射};
\end{axis}
\end{tikzpicture}
\caption{两状态概率单纯形上的平稳固定点：映射$r\mapsto0.2+0.5r$把任意初值逐步拉向$r_\star=0.4$}\label{fig:V5-C01-stationary-fixed-point}
\end{figure}

\clearpage
\typeout{SLFIG-HARNESS-END FIG-P556-01}

\typeout{SLFIG-HARNESS-BEGIN FIG-P556-02}
% v2.3.1 figure UID: FIG-P556-02
% Proposed post-v2.3.1 polish for FIG-P556-02: protect key-label size and stroke hierarchy.
\tikzset{
  slfig-FIG-P556-02/.style={font=\fontsize{9.2pt}{11.0pt}\selectfont,every path/.append style={line cap=round,line join=round}},
  slfig-FIG-P556-02-title/.style={font=\fontsize{10.4pt}{12.4pt}\selectfont\bfseries,anchor=north},
  slfig-FIG-P556-02-formula/.style={font=\fontsize{9.2pt}{11.0pt}\selectfont},
  slfig-FIG-P556-02-answer/.style={font=\fontsize{8.8pt}{10.5pt}\selectfont,text=SLGray,anchor=south},
}
\begin{figure}[htbp]
\centering
\begin{tikzpicture}[slfig-FIG-P556-02,>=Stealth,
  every node/.style={font=\fontsize{9.2pt}{11.0pt}\selectfont,align=center},
  card/.style={draw=SLRuleGray,rounded corners=2pt,fill=white,
    minimum width=41mm,minimum height=47mm,inner sep=5.5pt,line width=.55pt},
  state/.style={circle,draw=SLBlue,fill=SLBlue!7,minimum size=6.5mm,inner sep=0pt,
    font=\fontsize{9.4pt}{11.2pt}\selectfont,line width=.76pt},
  edge/.style={-{Stealth[length=1.5mm]},draw=SLTeal,line width=.82pt},
]
  \matrix (cards) [matrix of nodes,nodes={card},column sep=3.6mm] {
    |[name=irr]| {} & |[name=per]| {} & |[name=rec]| {} \\
  };

  \node[slfig-FIG-P556-02-title,text=SLBlue] at ([yshift=-3mm]irr.north) {不可约性};
  \node[state] (i1) at ([xshift=-9mm,yshift=7mm]irr.center) {$1$};
  \node[state] (i2) at ([xshift=9mm,yshift=7mm]irr.center) {$2$};
  \draw[edge] (i1) to[bend left=18] (i2);
  \draw[edge] (i2) to[bend left=18] (i1);
  \node[slfig-FIG-P556-02-formula] at ([yshift=-7mm]irr.center) {$i\leftrightarrow j$（支撑图连通）};
  \node[slfig-FIG-P556-02-answer] at ([yshift=3mm]irr.south)
    {回答：状态能否互相到达？};

  \node[slfig-FIG-P556-02-title,text=SLTeal] at ([yshift=-3mm]per.north) {周期性};
  \draw[draw=SLTextGray,line width=.60pt] ([xshift=-14mm,yshift=7mm]per.center)--([xshift=14mm,yshift=7mm]per.center);
  \foreach \x/\lab in {-12/2,-3/4,6/6,15/8}{
    \fill[SLTeal] ([xshift=\x mm,yshift=7mm]per.center) circle (1.4pt);
  }
  \node[slfig-FIG-P556-02-formula] at ([yshift=-7mm]per.center)
    {\shortstack{$d(i)=\gcd\{n\ge1:$\\[-.4mm]$P_{ii}^{(n)}>0\}$}};
  \node[slfig-FIG-P556-02-answer] at ([yshift=3mm]per.south)
    {回答：允许哪些回返步长？};

  \node[slfig-FIG-P556-02-title,text=SLGold] at ([yshift=-3mm]rec.north) {正常返};
  \node[state,draw=SLGold,fill=SLGold!8] (ri) at ([yshift=7mm]rec.center) {$i$};
  \draw[-{Stealth[length=1.6mm]},draw=SLGold,line width=.88pt]
    (ri) to[loop right,min distance=12mm] (ri);
  \node[slfig-FIG-P556-02-formula] at ([yshift=-7mm]rec.center) {$\mathbb E_i[\tau_i^+]<\infty$};
  \node[slfig-FIG-P556-02-answer] at ([yshift=3mm]rec.south)
    {回答：平均多久回到自身？};

  \node[draw=SLRuleGray,rounded corners=2pt,fill=SLSoftGray,
    minimum width=112mm,minimum height=7mm,inner xsep=4pt,line width=.58pt,
    font=\fontsize{9.2pt}{11.0pt}\selectfont] at ([yshift=-8mm]cards.south)
    {结构性质不可相互替代：连通性、回返节律与期望回返时间必须分别检查};
\end{tikzpicture}
\caption{三类结构必须分别判断：支撑图连通性决定不可约性，正回返时长的最大公约数决定周期性}\label{fig:V5-C01-chain-properties}
\end{figure}

\clearpage
\typeout{SLFIG-HARNESS-END FIG-P556-02}

\typeout{SLFIG-HARNESS-BEGIN FIG-P556-03}
% v2.3.1 figure UID: FIG-P556-03
% Proposed post-v2.3.1 polish for FIG-P556-03: protect key-label size and stroke hierarchy.
\tikzset{
  slfig-FIG-P556-03/.style={font=\fontsize{9.2pt}{11.0pt}\selectfont,every path/.append style={line cap=round,line join=round}},
  slfig-FIG-P556-03-title/.style={font=\fontsize{10.4pt}{12.4pt}\selectfont\bfseries,text=SLBlue},
  slfig-FIG-P556-03-verdict-text/.style={font=\fontsize{9.2pt}{11.0pt}\selectfont,anchor=west},
}
\begin{figure}[htbp]
\centering
\begin{tikzpicture}[slfig-FIG-P556-03,>=Stealth,
  every node/.style={font=\fontsize{9.2pt}{11.0pt}\selectfont,align=center},
  state/.style={circle,draw=SLBlue,fill=SLBlue!7,minimum size=9mm,inner sep=0pt,
    font=\fontsize{9.4pt}{11.2pt}\selectfont,line width=.82pt},
  flow/.style={-{Stealth[length=2mm]},draw=SLTeal,line width=1pt},
  lab/.style={font=\fontsize{9.2pt}{11.0pt}\selectfont,fill=white,inner sep=1.2pt},
  verdict/.style={draw=SLRuleGray,rounded corners=2pt,fill=white,
    minimum width=38mm,minimum height=14mm,line width=.56pt},
]
  \node[slfig-FIG-P556-03-title] at (0,2.15) {逐边双向概率流};
  \node[state] (x) at (-2.15,.85) {$x$};
  \node[state] (y) at (2.15,.85) {$y$};
  \draw[flow] (x) to[bend left=18]
    node[lab,above] {$\pi(x)K(x,y)$} (y);
  \draw[flow] (y) to[bend left=18]
    node[lab,below,yshift=2.4mm] {$\pi(y)K(y,x)$} (x);
  \node[draw=SLGold,rounded corners=2pt,fill=SLGold!8,
    minimum width=63mm,minimum height=8mm,line width=.72pt,
    font=\fontsize{9.4pt}{11.2pt}\selectfont] at (0,-.25)
    {$\pi(x)K(x,y)=\pi(y)K(y,x)$};

  \node[verdict] (v1) at (-4.45,-1.65) {};
  \node[verdict] (v2) at (0,-1.65) {};
  \node[verdict] (v3) at (4.45,-1.65) {};
  \node[circle,draw=SLTeal,text=SLTeal,minimum size=5.5mm,inner sep=0pt,
    line width=.72pt,font=\fontsize{9.2pt}{11.0pt}\selectfont]
    at ([xshift=5mm]v1.west) {$\checkmark$};
  \node[slfig-FIG-P556-03-verdict-text] at ([xshift=10mm]v1.west) {可推出平稳性};
  \node[circle,draw=SLRed,text=SLRed,minimum size=5.5mm,inner sep=0pt,
    line width=.72pt,font=\fontsize{9.2pt}{11.0pt}\selectfont]
    at ([xshift=5mm]v2.west) {$\times$};
  \node[slfig-FIG-P556-03-verdict-text] at ([xshift=10mm]v2.west) {不能推出连通};
  \node[circle,draw=SLRed,text=SLRed,minimum size=5.5mm,inner sep=0pt,
    line width=.72pt,font=\fontsize{9.2pt}{11.0pt}\selectfont]
    at ([xshift=5mm]v3.west) {$\times$};
  \node[slfig-FIG-P556-03-verdict-text] at ([xshift=10mm]v3.west) {不能推出唯一};

  % 极小断开可逆链反例：A=I_2。
  \node[state,minimum size=6.5mm,font=\fontsize{9.2pt}{11.0pt}\selectfont] (c1) at (-.65,-2.80) {$1$};
  \node[state,minimum size=6.5mm,font=\fontsize{9.2pt}{11.0pt}\selectfont] (c2) at (.65,-2.80) {$2$};
  \draw[flow,line width=.68pt] (c1) to[loop left,min distance=8mm] (c1);
  \draw[flow,line width=.68pt] (c2) to[loop right,min distance=8mm] (c2);
  \node[anchor=north west,yshift=-.8mm,
    font=\fontsize{8.8pt}{10.5pt}\selectfont,text=SLGray] at (1.15,-2.80)
    {$A=I_2$：可逆但断开，平稳分布不唯一};
\end{tikzpicture}
\caption{流量对称可证平稳，却不能证连通或唯一}\label{fig:V5-C01-detailed-balance}
\end{figure}

\clearpage
\typeout{SLFIG-HARNESS-END FIG-P556-03}

\typeout{SLFIG-HARNESS-BEGIN FIG-P558-01}
% v2.3.1 figure UID: FIG-P558-01
% Proposed post-v2.3.1 polish for FIG-P558-01: protect key-label size and stroke hierarchy.
\tikzset{slfig-FIG-P558-01/.style={every path/.append style={line cap=round,line join=round}}}
\pgfplotsset{slfig-FIG-P558-01-axis/.style={
  tick label style={font=\fontsize{8.6pt}{10pt}\selectfont},
  label style={font=\fontsize{9.2pt}{11pt}\selectfont},clip=false}}
\begin{figure}[htbp]
\centering
\begin{tikzpicture}[slfig-FIG-P558-01,>=Stealth,
  every node/.style={font=\fontsize{9.2pt}{11pt}\selectfont},
  state/.style={circle,draw=SLBlue,fill=SLBlue!7,minimum size=8mm,line width=.82pt,
    font=\fontsize{9.4pt}{11.2pt}\selectfont},
  edge/.style={-{Stealth[length=1.7mm]},draw=SLTeal,line width=.90pt},
  loop/.style={-{Stealth[length=1.7mm]},draw=SLGold,line width=.80pt},
  lab/.style={font=\fontsize{8.8pt}{10.4pt}\selectfont,fill=white,inner sep=1pt},
]
  \begin{scope}[shift={(-3.65,1.35)}]
    \node[font=\fontsize{10.4pt}{12.4pt}\selectfont\bfseries,text=SLBlue] at (0,1.58) {简单游走};
    \node[state] (a1) at (-1.75,0) {$1$};
    \node[state] (a2) at (0,0) {$2$};
    \node[state] (a3) at (1.75,0) {$3$};
    \draw[edge] (a1) to[bend left=18] node[lab,above] {$1$} (a2);
    \draw[edge] (a2) to[bend left=18] node[lab,below] {$1/2$} (a1);
    \draw[edge] (a2) to[bend left=18] node[lab,above] {$1/2$} (a3);
    \draw[edge] (a3) to[bend left=18] node[lab,below] {$1$} (a2);
  \end{scope}
  \begin{scope}[shift={(3.65,1.35)}]
    \node[font=\fontsize{10.4pt}{12.4pt}\selectfont\bfseries,text=SLTeal] at (0,1.58) {惰性游走};
    \node[state] (b1) at (-1.75,0) {$1$};
    \node[state] (b2) at (0,0) {$2$};
    \node[state] (b3) at (1.75,0) {$3$};
    \draw[edge] (b1) to[bend left=18] node[lab,above] {$1/2$} (b2);
    \draw[edge] (b2) to[bend left=18] node[lab,below] {$1/4$} (b1);
    \draw[edge] (b2) to[bend left=18] node[lab,above] {$1/4$} (b3);
    \draw[edge] (b3) to[bend left=18] node[lab,below] {$1/2$} (b2);
    \draw[loop] (b1) to[loop left,min distance=9mm] node[lab,left] {$1/2$} (b1);
    \draw[loop] (b2) to[loop above,min distance=7mm] node[lab,right] {$1/2$} (b2);
    \draw[loop] (b3) to[loop right,min distance=9mm] node[lab,right] {$1/2$} (b3);
  \end{scope}

  \begin{axis}[slfig-FIG-P558-01-axis,
    at={(-6.05cm,-.35cm)},anchor=north west,width=5.1cm,height=2.15cm,
    xmin=0,xmax=8,ymin=.18,ymax=.82,axis lines=left,
    xtick={0,2,4,6,8},ytick={.25,.5,.75},
    tick label style={font=\fontsize{8.6pt}{10pt}\selectfont},label style={font=\fontsize{9.2pt}{11pt}\selectfont},
    xlabel={$t$},ylabel={$P(X_t=2)$},clip=false]
    \addplot[draw=SLBlue,line width=.90pt,mark=*,mark size=1.35pt]
      coordinates {(0,.75)(1,.25)(2,.75)(3,.25)(4,.75)(5,.25)(6,.75)(7,.25)(8,.75)};
    \addplot[draw=SLGray,line width=.52pt,densely dashed] coordinates {(0,.5)(8,.5)};
    \node[anchor=south,yshift=.7mm,fill=white,inner sep=.8pt,
      font=\fontsize{8.8pt}{10.4pt}\selectfont,text=SLBlue] at (axis cs:4,.80) {奇偶振荡};
  \end{axis}
  \begin{axis}[slfig-FIG-P558-01-axis,
    at={(1.25cm,-.35cm)},anchor=north west,width=5.1cm,height=2.15cm,
    xmin=0,xmax=8,ymin=.18,ymax=.82,axis lines=left,
    xtick={0,2,4,6,8},ytick={.25,.5,.75},
    tick label style={font=\fontsize{8.6pt}{10pt}\selectfont},label style={font=\fontsize{9.2pt}{11pt}\selectfont},
    xlabel={$t$},ylabel={$P(X_t=2)$},clip=false]
    \addplot[draw=SLTeal,line width=.90pt,mark=*,mark size=1.35pt]
      coordinates {(0,.75)(1,.375)(2,.5625)(3,.46875)(4,.5156)(5,.4922)(6,.5039)(7,.4980)(8,.5010)};
    \addplot[draw=SLGray,line width=.52pt,densely dashed] coordinates {(0,.5)(8,.5)};
    \node[anchor=south,yshift=.7mm,fill=white,inner sep=.8pt,
      font=\fontsize{8.8pt}{10.4pt}\selectfont,text=SLTeal] at (axis cs:4,.80) {阻尼收敛};
  \end{axis}
  \node[draw=SLRuleGray,rounded corners=2pt,fill=SLSoftGray,line width=.55pt,
    minimum width=95mm,minimum height=7mm,align=center,
    font=\fontsize{9.2pt}{11pt}\selectfont] at (0,-3.15)
    {共享平稳分布 $\pi=(1/4,1/2,1/4)$；惰性自环把周期从 $2$ 降为 $1$};
\end{tikzpicture}
\caption{三节点路径的简单游走与惰性游走：加入自环保持平稳分布并消除周期振荡}\label{fig:V5-C01-random-walk}
\end{figure}

\clearpage
\typeout{SLFIG-HARNESS-END FIG-P558-01}

\typeout{SLFIG-HARNESS-BEGIN FIG-P570-01}
% v2.3.1 figure UID: FIG-P570-01
% Proposed post-v2.3.1 polish for FIG-P570-01: protect key-label size and stroke hierarchy.
\tikzset{slfig-FIG-P570-01/.style={font=\fontsize{9.2pt}{11.0pt}\selectfont,every path/.append style={line cap=round,line join=round}}}
\begin{figure}[htbp]
\centering
\begin{tikzpicture}[slfig-FIG-P570-01,>=Stealth,
  every node/.style={font=\fontsize{9.2pt}{11pt}\selectfont,align=center},
  input/.style={draw=SLRuleGray,rounded corners=2pt,fill=SLSoftGray,minimum width=30mm,minimum height=8mm,line width=.58pt},
  core/.style={draw=SLBlue,rounded corners=2pt,fill=SLBlue!7,minimum width=36mm,minimum height=9mm,line width=.90pt,
    font=\fontsize{9.6pt}{11.5pt}\selectfont},
  method/.style={draw=SLTeal,rounded corners=2pt,fill=white,
    minimum width=36mm,minimum height=13mm,inner xsep=1mm,line width=.78pt},
  dep/.style={-{Stealth[length=1.8mm]},draw=SLBlue,line width=.90pt},
]
  \node[input] (exp) at (-3.8,2.2) {期望 / 积分};
  \node[input] (iid) at (0,2.2) {独立抽样};
  \node[input] (law) at (3.8,2.2) {大数规律};
  \node[core] (mc) at (0,.75) {Monte Carlo 估计};
  \draw[dep] (exp)--(mc); \draw[dep] (iid)--(mc); \draw[dep] (law)--(mc);

  \node[method] (inv) at (-4.0,-.85) {\makebox[32mm][c]{\tikz\draw[SLBlue,line width=.68pt] (0,0)--(.35,.35);\quad 逆变换}};
  \node[method,densely dashed] (rej) at (0,-.85) {\makebox[32mm][c]{判定 $U\le r$\quad 接受--拒绝}};
  \node[method,double,double distance=1pt] (is) at (4.0,-.85) {\makebox[32mm][c]{$\vert\!\vert\!\vert$\quad 重要性抽样}};
  \draw[dep] (mc)--(inv); \draw[dep] (mc)--(rej); \draw[dep] (mc)--(is);

  \node[draw=SLGold,rounded corners=2pt,fill=SLGold!7,
    minimum width=116mm,minimum height=9mm,line width=.74pt,
    font=\fontsize{9.2pt}{11pt}\selectfont] (diag) at (0,-2.55)
    {共同出口：误差 / ESS / 支持覆盖诊断};
  \draw[dep] (inv)--(diag); \draw[dep] (rej)--(diag); \draw[dep] (is)--(diag);
  \node[font=\fontsize{8.6pt}{10.2pt}\selectfont,text=SLGray] at (0,-3.20)
    {三种同尺寸节点表示并列方法；边框/图标只编码方法身份，不表示先后依赖};
\end{tikzpicture}
\caption{本章知识依赖：由期望、积分与独立抽样进入蒙特卡罗估计，再分别学习三种直接采样或重加权方法，最后统一进行误差和支持诊断}\label{fig:V5-C02-dependency}
\end{figure}

\clearpage
\typeout{SLFIG-HARNESS-END FIG-P570-01}

\end{document}
